% =====================================================================
%  Chapter 19 — Which Origin Is True
% =====================================================================

\chapter{Which Origin Is True}
\label{ch:origin}

The second clause of this book's title asks where dark matter came from.
After five chapters of mechanisms, the honest answer is that the question is
not currently decidable, and that the reasons it is not decidable are
instructive enough to be worth stating precisely.

\section{What would make the question decidable}The origin question has three
possible resolutions, and they have different characters.

The first is \emph{discovery}. If a particle is detected, its properties
constrain the production mechanism: its mass, its couplings and its
abundance would determine which of the mechanisms in Chapter~\ref{ch:production}
is consistent, and in most cases would determine it uniquely. The QCD axion
is the model case: detection plus a measurement of the local density and the
velocity distribution would identify the cosmological scenario.

The second is \emph{comparative elimination}. If the cosmological history is
measured --- the inflationary power spectrum, the spectral distortions, the
expansion history --- the mechanisms that depend on it would be eliminated
along with it, and the survivors would be a much smaller set. This route
requires no detector and no new facility.

The third is \emph{logical simplification}. Some candidate combinations are
jointly excluded by existing data without any new measurement: a thermal
relic of a mass that is both excluded by direct detection and inconsistent
with the free-streaming bound is excluded. The eliminations of the last
twenty years have already removed a large part of the joint space, and the
book's own survey of the candidate landscape is largely a record of those
eliminations.

\section{Why preference cannot settle it}The field has a well-developed set of
preferences, and they are not bad reasons, but they are not decisive.

The WIMP was preferred because it predicted the abundance. The axion is
preferred because it solves a different problem while producing the right
abundance. Ultralight bosons are preferred because they explain the
small-scale structure. Primordial black holes are preferred because they
require no new physics.

Each of these is a reason to look first, not a reason to believe. A theory
preferred for aesthetic reasons has no protection, and the last two decades
have repeatedly demonstrated that the aesthetic preference and the
experimental outcome come apart. What would be illegitimate is to treat the
preference as if it were a prediction.

\section{The prior problem}The deepest issue is that the origins question is a
question about a prior, and there is no principled way to set a prior on the
early universe.

The density of dark matter is what it is, and every mechanism either
assumes it or explains it with a coincidence. The freeze-out mechanism
explains it with a coincidence that is remarkable and is itself unexplained.
The misalignment mechanism explains it with an initial angle. The
gravitational and Hawking mechanisms explain it with a primordial spectrum
that is measured elsewhere for other reasons and happens to have the right
amplitude at the right scale --- which is the most satisfying of the
available answers, because the spectrum is independently constrained.

The absence of a principled prior means that the field cannot converge on the
origins question by accumulating theoretical arguments. It can only converge
by measuring. This is the central argument of the book, and it is worth
stating in its strongest form: the origins question is not upstream of the
experimental programme, it is downstream of it.

\section{The honest position}There is no known origin of dark matter. The density
is known to two per cent; the mechanism that explained it with no free
parameter has been excluded; the mechanisms that remain either assume the
density or tie it to an independently measured spectrum; and no measurement
performed so far has discriminated between them.

This is not a failure of the field. It is a statement about the current
state of a young subject, and it has an actionable consequence: the
measurement that would settle the question most cheaply is a measurement of
the primordial spectrum and the CMB spectral distortions, and that is an
observatory experiment rather than a laboratory one.

\begin{keytakeaways}
\item The origins question is decidable by three routes: discovery of the
particle, comparative elimination through a measured cosmological history,
or joint elimination of inconsistent combinations.
\item Aesthetic preference is a reason to look first and not a reason to
believe, and the last two decades have separated the two repeatedly.
\item There is no principled prior on the early universe, so the field cannot
converge on the origins question by argument; it converges by measurement.
\item The cheapest decisive test is a measurement of the primordial power
spectrum and the CMB spectral distortions.
\end{keytakeaways}
